% GENERATED FILE. Edit canonical lessons, then run npm run generate:books.
% canonical-source-hash: fnv1a64:706b3637f8d66932
% canonical-lessons: JA-C174-nabe, JA-C174-houchou, JA-C174-kizamu, JA-C174-osara, JA-C174-moru

\chapter{Pots, Knives and Plates}
\label{ch:ja-pots-knives-and-plates}
\hlchaptermodality{174}

\begin{chapteropening}
\textbf{By the end of this chapter:} \emph{I can name kitchen tools and say what I do with them.}

Complete the last lesson of chapter 174: I can name kitchen tools and say what I do with them.
\end{chapteropening}

\section[\texorpdfstring{\ja{なべ} (nabe)}{なべ (nabe)}]{\ja{なべ} (nabe) — pot; saucepan, hot pot}

\label{lesson:JA-C174-nabe}



\begin{quote}
Before the new one: say the Japanese for to be left over, then the Japanese for to savour.
\end{quote}

\subsection*{You'll want to know: \ja{なべ}}
\textbf{\ja{なべ}} — \emph{nabe} — \textquotedblleft{}pot; saucepan, hot pot\textquotedblright{}.

\begin{quote}
\textbf{\ja{ふゆは} \ja{なべを} \ja{よく} \ja{たべます。}} — \emph{fuyu wa nabe o yoku tabemasu} — \textquotedblleft{}In winter we often eat hot pot.\textquotedblright{}
\end{quote}

\begin{cousinweb}
\textbf{\ja{なべ}} (\emph{nabe}) is a native Japanese word (wago), inherited rather than borrowed from Chinese; its kanji was attached later as a kun'yomi; its older origin is not settled (JMdict, KANJIDIC2).
\end{cousinweb}

\begin{grammarlens}[title={What you've built}]
One more word for everyday life at A2.
\end{grammarlens}

\subsection*{Your turn}
\begin{itemize}
\raggedright
  \item \emph{Say it:} \emph{nabe}
  \item \emph{Say it:} \emph{nabe}, once more
  \item \emph{Say it:} \emph{fuyu wa nabe o yoku tabemasu} — In winter we often eat hot pot.
  \item \emph{Recall:} say the Japanese for to be left over, then the Japanese for to savour, then say \emph{nabe} again
\end{itemize}

\begin{tcolorbox}[breakable,colback=teal!4,colframe=teal!35!black,title={Before you move on}]
What does \ja{なべ} mean? (\textquotedblleft{}Pot; saucepan, hot pot\textquotedblright{}.) Say it once more.
\end{tcolorbox}
\section[\texorpdfstring{\ja{ほうちょう} (hōchō)}{ほうちょう (hōchō)}]{\ja{ほうちょう} (hōchō) — kitchen knife}

\label{lesson:JA-C174-houchou}



\begin{quote}
Before the new one: say the Japanese for to savour, then the Japanese for pot.
\end{quote}

\subsection*{You'll want to know: \ja{ほうちょう}}
\textbf{\ja{ほうちょう}} — \emph{hōchō} — \textquotedblleft{}kitchen knife\textquotedblright{}.

\begin{quote}
\textbf{\ja{ちちは} \ja{ほうちょうを} \ja{よく} \ja{つかいます。}} — \emph{chichi wa hōchō o yoku tsukaimasu} — \textquotedblleft{}My father often uses a kitchen knife.\textquotedblright{}
\end{quote}

\begin{cousinweb}
\textbf{\ja{ほうちょう}} (\emph{hōchō}) is Sino-Japanese (kango): two Chinese characters read with their on'yomi, readings borrowed from Chinese. Today they are \emph{hō}, \textquotedblleft{}wrap\textquotedblright{}, and \emph{chō}, a character that also counts tools (Middle Chinese \emph{teng}); JMdict also lists an older spelling whose first character the Guangyun glosses \textquotedblleft{}kitchen\textquotedblright{} (\emph{baew}), so the history of the pair is not settled (JMdict, KANJIDIC2; Middle Chinese in Baxter's transcription).
\end{cousinweb}

\begin{grammarlens}[title={What you've built}]
One more word for everyday life at A2.
\end{grammarlens}

\subsection*{Your turn}
\begin{itemize}
\raggedright
  \item \emph{Say it:} \emph{hōchō}
  \item \emph{Say it:} \emph{hōchō}, once more
  \item \emph{Say it:} \emph{chichi wa hōchō o yoku tsukaimasu} — My father often uses a kitchen knife.
  \item \emph{Recall:} say the Japanese for to savour, then the Japanese for pot, then say \emph{hōchō} again
\end{itemize}

\begin{tcolorbox}[breakable,colback=teal!4,colframe=teal!35!black,title={Before you move on}]
What does \ja{ほうちょう} mean? (\textquotedblleft{}Kitchen knife\textquotedblright{}.) Say it once more.
\end{tcolorbox}
\section[\texorpdfstring{\ja{きざむ} (kizamu)}{きざむ (kizamu)}]{\ja{きざむ} (kizamu) — to chop; to mince, to carve}

\label{lesson:JA-C174-kizamu}



\begin{quote}
Before the new one: say the Japanese for pot, then the Japanese for kitchen knife.
\end{quote}

\subsection*{You'll want to know: \ja{きざむ}}
\textbf{\ja{きざむ}} — \emph{kizamu} — \textquotedblleft{}to chop; to mince, to carve\textquotedblright{}.

\begin{quote}
\textbf{\ja{ほうちょうで} \ja{だいこんを} \ja{きざみます。}} — \emph{hōchō de daikon o kizamimasu} — \textquotedblleft{}I chop a daikon with a kitchen knife.\textquotedblright{}
\end{quote}

Polite \textbf{\ja{きざみます}}, \emph{kizamimasu}; a godan verb.

\begin{cousinweb}
\textbf{\ja{きざむ}} (\emph{kizamu}) is a native Japanese word (wago), inherited rather than borrowed from Chinese; its kanji was attached later as a kun'yomi (JMdict, KANJIDIC2).
\end{cousinweb}

\begin{grammarlens}[title={What you've built}]
One more word for everyday life at A2.
\end{grammarlens}

\subsection*{Your turn}
\begin{itemize}
\raggedright
  \item \emph{Say it:} \emph{kizamu}
  \item \emph{Say it:} \emph{kizamu}, once more
  \item \emph{Say it:} \emph{hōchō de daikon o kizamimasu} — I chop a daikon with a kitchen knife.
  \item \emph{Recall:} say the Japanese for pot, then the Japanese for kitchen knife, then say \emph{kizamu} again
\end{itemize}

\begin{tcolorbox}[breakable,colback=teal!4,colframe=teal!35!black,title={Before you move on}]
What does \ja{きざむ} mean? (\textquotedblleft{}To chop; to mince, to carve\textquotedblright{}.) Say it once more.
\end{tcolorbox}
\section[\texorpdfstring{\ja{おさら} (osara)}{おさら (osara)}]{\ja{おさら} (osara) — plate; dish}

\label{lesson:JA-C174-osara}



\begin{quote}
Before the new one: say the Japanese for kitchen knife, then the Japanese for to chop.
\end{quote}

\subsection*{You'll want to know: \ja{おさら}}
\textbf{\ja{おさら}} — \emph{osara} — \textquotedblleft{}plate; dish\textquotedblright{}.

\begin{quote}
\textbf{\ja{さかなを} \ja{おさらに} \ja{だします。}} — \emph{sakana o osara ni dashimasu} — \textquotedblleft{}I serve the fish on a plate.\textquotedblright{}
\end{quote}

\begin{cousinweb}
\textbf{\ja{おさら}} (\emph{osara}) is built from polite \emph{o-} and native \emph{sara}, \textquotedblleft{}dish\textquotedblright{}, inherited rather than borrowed; its kanji was attached later as a kun'yomi (JMdict, KANJIDIC2).
\end{cousinweb}

\begin{grammarlens}[title={What you've built}]
One more word for everyday life at A2.
\end{grammarlens}

\subsection*{Your turn}
\begin{itemize}
\raggedright
  \item \emph{Say it:} \emph{osara}
  \item \emph{Say it:} \emph{osara}, once more
  \item \emph{Say it:} \emph{sakana o osara ni dashimasu} — I serve the fish on a plate.
  \item \emph{Recall:} say the Japanese for kitchen knife, then the Japanese for to chop, then say \emph{osara} again
\end{itemize}

\begin{tcolorbox}[breakable,colback=teal!4,colframe=teal!35!black,title={Before you move on}]
What does \ja{おさら} mean? (\textquotedblleft{}Plate; dish\textquotedblright{}.) Say it once more.
\end{tcolorbox}
\section[\texorpdfstring{\ja{もる} (moru)}{もる (moru)}]{\ja{もる} (moru) — to dish up; to serve (food), to heap}

\label{lesson:JA-C174-moru}



\begin{quote}
Before the new one: say the Japanese for to chop, then the Japanese for plate.
\end{quote}

\subsection*{You'll want to know: \ja{もる}}
\textbf{\ja{もる}} — \emph{moru} — \textquotedblleft{}to dish up; to serve (food), to heap\textquotedblright{}.

\begin{quote}
\textbf{\ja{おさらに} \ja{ごはんを} \ja{もります。}} — \emph{osara ni gohan o morimasu} — \textquotedblleft{}I dish up rice onto a plate.\textquotedblright{}
\end{quote}

Polite \textbf{\ja{もります}}, \emph{morimasu}; a godan verb.

\begin{cousinweb}
\textbf{\ja{もる}} (\emph{moru}) is a native Japanese word (wago), inherited rather than borrowed from Chinese; its kanji was attached later as a kun'yomi (JMdict, KANJIDIC2).
\end{cousinweb}

\begin{grammarlens}[title={What you've built}]
One more word for everyday life at A2.
\end{grammarlens}

\subsection*{Your turn}
\begin{itemize}
\raggedright
  \item \emph{Say it:} \emph{moru}
  \item \emph{Say it:} \emph{moru}, once more
  \item \emph{Say it:} \emph{osara ni gohan o morimasu} — I dish up rice onto a plate.
  \item \emph{Recall:} say the Japanese for to chop, then the Japanese for plate, then say \emph{moru} again
\end{itemize}

\begin{tcolorbox}[breakable,colback=teal!4,colframe=teal!35!black,title={Before you move on}]
What does \ja{もる} mean? (\textquotedblleft{}To dish up; to serve (food), to heap\textquotedblright{}.) Say it once more.
\end{tcolorbox}
